\chapter{Introduction}

\section{What the system is}

\vapor{} has two parts that share one discipline.

The first is a \textbf{compiler and runtime for tensor programs}. A program is a term of a small free
algebra (element-wise operations, contractions, reductions, the operators of modern neural networks)
whose meaning is fixed by an exact reference implementation, the \emph{oracle}. The compiler lowers a
program to a portable kernel IR, allocates registers without spilling, and encodes machine code for
four instruction sets and for Vulkan compute shaders. The runtime executes that code in a separate,
sandboxed process, keeps model weights in content-addressed shared memory, serves models through an
OpenAI-compatible interface and records a receipt for every answer.

The second is a \textbf{laboratory}: deciders, languages, search, numerical workbenches, statistical
tools, document readers and desks for engineering and finance, all reachable from one command-line
interpreter, a browser console, an MCP server for agents, a language server for editors, and a small
editor of its own.

What the two share is the rule that a result is never just a value. A compiled kernel comes with the
evidence that it computes the oracle's bits; an \emph{optimal} integer program comes with a
branch-and-bound tree that separate code checks; a refuted conservation law comes with the point that
refutes it; a statistical claim comes with a control that a broken implementation would have passed.

\section{Three properties}

Three properties run through every chapter.

\paragraph{Determinism.} Under the \emph{canonical} numerical policy, a program produces the same bits
on every substrate: x86-64 with AVX2 or AVX-512, AArch64 with NEON, RISC-V with RVV (native or
interpreted), Vulkan (on any driver) and the oracle. The bits also do not depend on the number of
threads, on how many sequences share a decoding step, on where the key--value cache lives, or on which
replica serves a request. Parallelism is drawn from independent outputs, never from reassociating a sum
(Chapter~\ref{ch:semantics}).

\paragraph{Certification.} Wherever a claim can be checked by a procedure simpler than the one that
produced it, the system ships the certificate and the checker. Register allocations are revalidated by
a checker proved sound in Lean~4. Kernels climb a ladder of differential checks that ends in a signed,
reproducible certificate. Deciders return witnesses, Farkas vectors, DRUP proofs or hedges, and each is
checked by code that shares nothing with the search (Chapters~\ref{ch:verification}
and~\ref{ch:deciders}).

\paragraph{Containment.} Nothing generated runs inside the BEAM: there are no NIFs. Machine code runs in
a static, libc-free worker under seccomp, with W\^{}X pages and a watchdog. Work on untrusted input runs
under one containment primitive, the \emph{hermetic seal}, which caps memory (on-heap and off-heap) and
time. Models enter only through the model airlock; bytes from the network enter only through the
network airlock, and only a person can open it (Chapters~\ref{ch:models} and~\ref{ch:network}).

\section{How to read this document}

Chapter~\ref{ch:overview} gives the layers and the trust boundaries. Chapters~\ref{ch:semantics}
to~\ref{ch:verification} describe the core: the algebra and its numerical semantics, the compiler, the
runtime and the verification ladder. Chapters~\ref{ch:models} and~\ref{ch:network} describe the two
airlocks. Chapters~\ref{ch:deciders} and~\ref{ch:languages} describe the certified deciders and the
languages built on them. Chapter~\ref{ch:evidence} describes how quality is measured and how the test
suite is organised, Chapter~\ref{ch:interfaces} the interfaces, and Chapter~\ref{ch:limits} what the
system does not do. The appendix maps components to modules.

Every statement of fact in this document corresponds to a test or a check in the repository, named in
the system's own documents (\texttt{docs/}); this monograph gives the architecture and leaves the
contestation commands to them.
